% =========================================================
% III. PROBLEM FORMULATION
% =========================================================
\section{Problem Formulation}
\label{sec:problem_formulation}

We consider multi-hop legal question answering, where a conclusion may require
connecting several statutory conditions, intermediate events, and legal
consequences. Let the input question be $q\in\mathcal{Q}$. Legal knowledge is
represented by a directed graph $G=(V,E)$ and an associated rule set
$\mathcal{R}$. Each node $v_i\in V$ denotes a normalized legal event, while an
edge summarizes one or more statutory condition--consequence rules. The rule
set, rather than graph reachability alone, defines the structural inference
semantics.

Given $q$, the system retrieves relevant events and constructs candidate
reasoning paths
\begin{equation}
\mathcal{P}(q)=\{P_1,P_2,\ldots,P_K\},
\end{equation}
where
\begin{equation}
P_i=v_0\xrightarrow{r_1}v_1\xrightarrow{r_2}\cdots
\xrightarrow{r_h}v_h.
\end{equation}
The primary path is selected by
\begin{equation}
P^*=\arg\max_{P_i\in\mathcal{P}(q)}S(P_i,q),
\end{equation}
where $S(P_i,q)$ combines semantic and structural relevance.

\subsection{Three-Valued Rule Semantics}

Each event is associated with a structural variable
\begin{equation}
X_i\in\{1,0,\bot\},
\end{equation}
where $1$, $0$, and $\bot$ denote \textsc{true}, \textsc{false}, and
\textsc{unknown}. The third value represents evidence that is not established
or is conflicting; it is not interpreted as a numeric value between zero and
one. This design follows three-valued fixed-point semantics for logic programs
\cite{kleene1952metamathematics,fitting1985kripke} and their linear-algebraic
extensions \cite{sato2020threevalued}.

To encode the state without assigning a numeric value to $\bot$, we use two
support channels
\begin{equation}
t_i=\mathbb{I}[X_i=1],\qquad
f_i=\mathbb{I}[X_i=0].
\label{eq:two_channel_state}
\end{equation}
Thus, $X_i=\bot$ gives $(t_i,f_i)=(0,0)$. Conflicting positive and negative
support is conservatively projected to $\bot$ by the runtime inference engine.

Let $\mathbf{B}^{+},\mathbf{B}^{-}\in\{0,1\}^{|\mathcal{R}|\times |V|}$ be
body-incidence matrices for literals requiring true and false states. For rule
$r$, let $k_r=\sum_i(B^{+}_{ri}+B^{-}_{ri})$. Its conjunctive body is satisfied
iff every literal has the required settled state:
\begin{equation}
b_r(\mathbf{t},\mathbf{f})=
\mathbb{I}\!\left[
\sum_i B^{+}_{ri}t_i+
\sum_i B^{-}_{ri}f_i=k_r
\right].
\label{eq:conjunctive_body}
\end{equation}
If $e_r$ indicates that at least one encoded exception is satisfied, the rule
activation is $a_r=b_r(1-e_r)$. Let $\mathbf{H}^{+}$ and $\mathbf{H}^{-}$ be
head-incidence matrices. Positive support for event $i$ is aggregated as
\begin{equation}
d_i^{+}=\mathbb{I}\!\left[\sum_r H^{+}_{ir}a_r\geq 1\right],
\label{eq:disjunctive_support}
\end{equation}
with an analogous expression for negative support. Consequently, literals
inside one rule are conjunctive, whereas independently activated rules with the
same head are alternative, disjunctive mechanisms. For example,
$(A\land B)\rightarrow Y$ is one rule with two body entries, while
$(A\lor B)\rightarrow Y$ is normalized into $A\rightarrow Y$ and
$B\rightarrow Y$. Disjunctive heads are outside the current rule schema.
This body/head incidence view is the sparse matrix counterpart of
linear-algebraic and tensor-space representations of definite, disjunctive,
and normal logic programs \cite{sakama2021tensor}.

The factual state is the fixed point
\begin{equation}
\mathbf{x}^{F}
=
\operatorname{Fix}(\Phi_{\mathcal{R}},\mathbf{x}^{(0)}),
\label{eq:factual_fixed_point}
\end{equation}
where $\Phi_{\mathcal{R}}$ activates the encoded rule mechanisms. The matrices
above define a formal algebraic representation. The implementation evaluates
the corresponding sparse rules symbolically and iterates to a fixed point; it
does not perform dense matrix multiplication. The engine supports multiple
body literals, but the legacy rule collection used in the reported experiment
maps each record to one positive antecedent and one positive consequent.
Therefore, conjunction is an engine-level capability and is not empirically
validated by the present benchmark.

\subsection{Topological Ablation and Structural Intervention}

Let $M(P^*)=\{v_1,\ldots,v_{h-1}\}$ denote the intermediate events of the
primary path. We distinguish two non-equivalent operations. First, a
node-deletion diagnostic forms the induced subgraph
\begin{equation}
G^{-M}=G[V\setminus\{M\}]
\end{equation}
and checks bounded reachability
\begin{equation}
\rho_{\mathrm{top}}(s,Y\mid M)
=
\mathbb{I}[Y\text{ is reachable from }s\text{ in }G^{-M}].
\label{eq:topological_ablation}
\end{equation}
This operation reports whether an alternative graph path is found within the
search limits. It is used only as a topology diagnostic or fallback and is not
the LegalSCM counterfactual definition.

Second, for an explicit mediator-removal query, LegalSCM applies the hard
intervention
\begin{equation}
do(X_M=0),
\label{eq:problem_do}
\end{equation}
where $0\equiv\mathrm{FALSE}$. The variable and the remaining rule structure
are preserved; every rule whose head is $X_M$ is disabled and the mechanism for
$X_M$ is replaced by the constant assignment $X_M:=0$. Equivalently,
\begin{equation}
\mathcal{R}^{do(X_M=0)}
=
\{r\in\mathcal{R}\mid \operatorname{head}(r)\neq X_M\}.
\end{equation}
The counterfactual world is recomputed as
\begin{equation}
\mathbf{x}^{CF}_{M}
=
\operatorname{Fix}\!\left(
\Phi_{\mathcal{R}^{do(X_M=0)}},
\mathbf{x}^{(0)}\cup\{X_M=0\}
\right).
\label{eq:cf_fixed_point}
\end{equation}
For a target $Y$, $M$ is necessary relative to the encoded normative model when
$x_Y^{F}=1$ and $x_{Y,M}^{CF}=0$; it is non-necessary when
$x_{Y,M}^{CF}=1$, and remains indeterminate when either world is unknown.
This is a hard intervention over the encoded rule model, without an exogenous
abduction layer, and must not be equated with deleting $M$ from $G$.

Finally, the system predicts a query-relative decision $z\in\mathcal{Z}$,
where
\begin{equation}
\begin{aligned}
\mathcal{Z}=\{&\textsc{supported},\textsc{uncertain},\\
&\textsc{reject-direct-claim}\},
\end{aligned}
\end{equation}
and generates an answer $a$ with statutory citations $C$:
\begin{equation}
f:(q,G,\mathcal{R})\rightarrow(P^*,z,\mathcal{E}^*,a,C),
\end{equation}
where $\mathcal{E}^*$ is the refined evidence supplied to the generation model.
The decision labels are distinct from both the event-state domain
$\{1,0,\bot\}$ and internal mediator statuses such as necessary or
non-necessary.
